The analysis presented in this work is based on a comprehensive set of local distance measurements and their corresponding calibrators, as detailed in Section~\ref{sec:meth} and conceptually shown in Fig.\,\ref{fig:Metro}. The dataset includes a broad array of distance indicators, spanning geometric anchors, calibrators, and tracers. These are combined using a statistically rigorous framework that accounts for shared systematics and covariances among measurements.
To maintain clarity and readability in the main text, we provide here only a brief summary of the datasets used to construct the \DN, summarized also in Table~\ref{tb:datasets}. Methods employed in the baseline are marked in bold. A more detailed description of the individual measurements, sources, and associated assumptions is given in Appendix~\ref{sec:data_appendix}. We employ peculiar velocity corrections as described in Appendix~\ref{app:sec:pec_vel}.

%\nilsc{Honestly, I think this section should be much more explanative about HOW these methods work, with just a tiny bit of info on which we use precisely, with the big info on WHAT data we use mostly already found in App. A. For example, for the JABG we say nothing of how it works, just which PI's have observed it in which programs, which IMHO is irrelevant information for any non-expert.}

% \riacomm{In keeping with the below listing methods (not objects), I have slightly rearranged and relabeled the first three paragraphs} \nilsc{I found the previous labeling more intuitive, since it separated out rungs.}

\paragraph{\textbf{Parallaxes}}
%Trigonometric parallaxes of Milky Way (MW) stars}}
We adopted trigonometric wide-angle parallaxes (henceforth: parallaxes) of Milky Way (MW) stars from the early third data release of the ESA {\it Gaia} mission \citep[GEDR3;][]{2016A&A...595A...1G,2021A&A...649A...1G}. GEDR3 parallaxes were corrected for systematics following \citet{2021A&A...649A...4L} and for additional residual offsets as determined for field Cepheids \citep{2021ApJ...908L...6R}. GEDR3 parallaxes of open stars clusters were determined using non-variable member stars \citep{2022ApJ...938...36R,2023A&A...672A..85C}. An additional seven narrow-angle parallaxes of Cepheids based on the HST/WFC3 spatial scanning technique were included \citep{2018ApJ...855..136R}. 
% \citep[7 Cepheids]{} 
% We adopt the trigonometric parallaxes of individual Cepheids within the Milky way, taking into account its extended nature (as opposed to other anchors, whose extent relative to their distance is often negligble). The distances are measured with high signal-to-noise ratio by the ESA {\it Gaia} mission~\citep{2016A&A...595A...1G}.

\paragraph{\textbf{Detached Eclipsing Binary (DEB) distances}} % to the Magellanic Clouds}}
The distances to both Magellanic Clouds have been determined using helium burning red giants in detached eclipsing binary systems (henceforth: DEBs). These distances are determined geometrically as the ratio of the physical stellar radii determined from full orbital solutions to the angular diameters determined by surface-brightness-color relations calibrated empirically using long-baseline interferometry. 
% solved orbits \nils{. This method uses the fact that the orbits closely determine the (eclipsed) star's linear radius to extreme precision. Estimating then the angular size of the star from the measured flux as well as the surface brightness (calibrated using an emprical surface-brightness-color relation) allows the determination of the distance}. 
% For the present work, we assume a 
The distance to the LMC from \cite{2019Natur.567..200P} contributes to the Cepheid calibration in the baseline analysis. An analogous distance to the SMC from \cite{2020ApJ...904...13G} was considered as part of a variant. 
% In our analysis, both galaxies are used as geometric anchors to calibrate the period-luminosity relation of classical Cepheids. 


\paragraph{\textbf{Megamaser distances}}
We adopted geometric distance measurements to the anchor galaxy NGC$\,$4258 and five additional galaxies out to the Hubble flow. The distances were determined via very long-baseline interferometric (VLBI) radio-wavelength observations of water megamasers in accretion disks surrounding central SMBHs of their host galaxies. Distances were determined by comparing the physical scale of the megamasers determined from a Keplerian disk model to the angular extent of the features \citep{2015IAUGA..2255730B,2019ApJ...886L..27R,2020ApJ...892...18K}. 

%% merged the following two together since it's the same method
% \paragraph{\textbf{NGC$\,$4258}}
% We adopt the geometric distance measurement to NGC$\,$4258 from~\citet{2019ApJ...886L..27R}, who used VLBI observations of water masers in the accretion disk around the central SMBH to model their orbital motion and determine the distance to the system.  In our analysis, NGC$\,$4258 is used as an anchor to calibrate the Cepheid variable, TRGB, JAGB, and Mira variable primary distance indicators.
% \paragraph{\textbf{Megamaser distances}}
% We use a sample of 5 megamaser-hosting galaxies as single-step tracers in the Hubble flow.  The distances to these galaxies were compiled by the Megamaser Cosmology Project  \citep[][MCP]{2015IAUGA..2255730B,2020ApJ...892...18K}, which used spectral monitoring and VLBI mapping observations to measure the orbital motions of the masers and thereby determine geometric distance measurements to each system.


\paragraph{\textbf{Cepheids}}
The Leavitt Law, or Period-Luminosity relation of classical Cepheids (throughout this work: Cepheids) is a well-understood consequence of the dependence of acoustic oscillations on stellar density, the relation between mass and luminosity, and the Stefan-Boltzmann law. The periods of their observed brightness variations link directly to their intrinsic luminosity  \citep{1912HarCi.173....1L} with particularly small scatter in the infrared, and this is well described by stellar evolution models \citep[e.g.,][]{2025arXiv250522512K}. Distances to Cepheids pulsating in the fundamental mode were measured using the reddening-free near-infrared Wesenheit magnitudes based on HST photometry published by \citet{2018ApJ...855..136R, 2021ApJ...908L...6R, 2022ApJ...938...36R, 2019ApJ...876...85R, 2022ApJ...934L...7R, 2022ApJ...940...64Y, 2024ApJ...973...30B}, and more recently with JWST by \citet{2023ApJ...956L..18R, 2024ApJ...977..120R, 2025arXiv250901667R}. The Cepheid Period-Luminosity relation was calibrated in anchor galaxies using geometric distances in the Milky Way, Magellanic Clouds, and NGC$\,$4258.

\paragraph{\textbf{Tip of the Red Giant Branch (TRGB)}}
The TRGB represents a recognizable feature in the color-magnitude diagram of galaxies that is caused by the nearly constant luminosity of the Helium flash of first-ascent red giant stars, which continue along their evolution towards the Horizontal Branch.
%sudden drop in luminosity once Helium burning starts, whose intrinsic luminosity is universal.}
The TRGB distances were measured using HST and JWST observations of resolved stars in nearby (D$<$30 Mpc) galaxies. These distances were measured by the Carnegie-Chicago Hubble Program (CCHP; \citealt{2019ApJ...882...34F, 2021ApJ...915...34H, 2025arXiv250311769H}), the Extragalactic Distance Database (EDD; \citealt{2009AJ....138..323T, 2021AJ....162...80A}), the SH0ES team \citep{2024ApJ...966...89A,2024ApJ...976..177L}, and the TRGB-SBF project team \citep{2024ApJ...973...83A, 2025ApJ...982...26A}. Each team used somewhat different reduction, analysis, and calibration techniques, but we find that the resulting distances are generally consistent across groups, as shown in Section~\ref{sec:results}.

\paragraph{J-region Asymptotic Giant Branch (JAGB)} The JAGB is a recognizable feature in the near-infrared color-magnitude diagram of galaxies attributed to post-third dredge-up carbon-rich intermediate-mass AGB stars. The JAGB distances used here are derived from data taken with {JWST} NIRCam as published by \citep[][SH0ES]{2024ApJ...966...20L,2025arXiv250205259L} and \citep[][CCHP]{2025ApJ...985..203F}; see also Appendix~\ref{sec:data_appendix}. In its current implementation, the JAGB method assumes a homogeneous stellar population with a constant average luminosity. Systematic uncertainties, e.g., asymmetric luminosity functions or metallicity effects, were considered following  \citet{2024ApJ...966...20L} and remain a subject of research \citep[e.g.,][]{2021ApJ...916...19Z,2024AA...691A.350M,2025ApJ...985..182L}.
% The JAGB method currently assumes a homogeneous stellar population in the J-region; while this may not be the case \nils{and there could be some} inhomogeneity, which can manifest as asymmetry of the J-region luminosity function, \nils{it} is accounted for as a systematic uncertainty \nils{determined} via measurement variants, for example as described in \cite{2024ApJ...966...20L}.

\paragraph{Mira distances}
Miras are high-amplitude long-period variable asymptotic giant branch stars that obey period-luminosity relations.
We use a sample of 3 Mira host galaxies consisting of the geometric anchor NGC$\,$4258 \citep{2018ApJ...857...67H} and two SNe~Ia calibrator galaxies, NGC$\,$1559 and M101 \citep{2020ApJ...889....5H, 2024ApJ...963...83H}, hosts of SN~2005df and SN~2011fe, respectively. All observations consisted of at least 10 epochs of HST WFC3/IR time series photometry spanning a minimum of one year, and distances were obtained using the Mira period-luminosity relation in the {F160W} bandpass. 


\paragraph{\textbf{Type Ia Supernovae (SNe~Ia)}}
SNe~Ia are standardizable candles whose light curve information (e.g., the duration) can be related to their intrinsic magnitude.
We incorporate five different SNe~Ia datasets, ranging from optical to near-infrared (NIR) wavelengths and spanning different SNe~Ia modeling methodologies (Spectral template PCA with SALT2, SALT3 and SNooPy v2.7, and template-free with BayesSN). Four of these datasets are used in the Hubble Flow: (1) Pantheon+ (optical) with SALT2 \citep{2022ApJ...938..113S,2022ApJ...938..111B}, (2) Carnegie Supernova Project (CSP) I \& II (with SNooPy v2.7; \citealt{2011AJ....141...19B,2024ApJ...970...72U}), (3) template-independent distances in NIR \citep{2023AA...679A..95G}, (4) optical+NIR samples processed with BayesSN \citep{2023MNRAS.524..235D}. A fifth dataset of 13 SNe~Ia is used to calibrate the distance to the Coma cluster \citep{2025ApJ...979L...9S}. % \riacomm{inconsistency: \ref{sec:eq_special:FP} and FP below mention 13 SNe Ia, one of which is used twice}

\paragraph{Type II Supernovae (SNe~II)}
Several methods exist to standardize SN~II magnitudes and derive distances (see \citealt{2024hct..book..177D}). In particular, the Standard Candle Method (SCM; \citealt{2002ApJ...566L..63H}) standardizes SN~II luminosities based on correlations with photospheric velocity (from H$_\beta$) and color.
We use a sample of 89 SNe~II in the Hubble flow at $z > 0.01$ as well as 14 calibrator supernovae from \citet{2020MNRAS.495.4860D,2022MNRAS.514.4620D}, compiled from CSP-I, LOSS, SDSS-II, SNLS, DES-SN, and SSP-HSC \citep{2017ApJ...835..166D,2017MNRAS.472.4233D,2020MNRAS.495.4860D}.

\paragraph{Expanding photosphere method (EPM) of Type IIP supernovae}
We adopted distances of Type IIP supernovae (SNe~IIP) determined by astrophysical modeling via the tailored expanding photosphere method (EPM) from \citet{2025AA...702A..41V}. Such distances relate the angular extent of SNe~IIP measured from light curves to their physical expansion determined by radiative transfer modeling of optical spectra \citep{2020A&A...633A..88V,2023AA...672A.129C}. The EPM can be applied to SNe~IIP in the Hubble flow without requiring additional calibration, although it depends on the accuracy of the underlying astrophysical modeling.


\paragraph{\textbf{Surface Brightness Fluctuations (SBF)}} Spatial fluctuations in the surface brightness of an otherwise smooth galaxy arise from the statistics of the discrete number of stars per pixel. The amplitude of these fluctuations depends inversely on the distance of the galaxy, as well as on age, metallicity, and other properties of the stellar population.
For galaxies with evolved stellar populations, the SBF method can be calibrated through an empirical relation between the intrinsic magnitude of the fluctuations and the galaxy's integrated color, used as a proxy for the stellar population properties. SBF measurements in the near-infrared are particularly useful because the fluctuations are bright at these wavelengths and can be well calibrated using optical colors.
For this study, we use a sample of 61 HST WFC3/IR SBF elliptical galaxy distances reaching out to 100 Mpc \citep{2025ApJ...987...87J,2021ApJS..255...21J,2021ApJ...911...65B}; distances to the 14 calibrators in the Virgo and Fornax galaxy clusters are based on TRGB; see Appendix~\ref{sec:eq_groups}.

\paragraph{Fundamental Plane (FP)}
The distance to an elliptical galaxy can be estimated by means of the fundamental plane relationship between central velocity dispersion, effective radius, and surface brightness.
Fundamental Plane distances were measured using the Dark Energy Spectroscopic Instrument (DESI) Early Data Release, analyzing 4191 early-type galaxies within $0.01 < z < 0.1$ with photometry from the DESI Legacy Imaging Surveys and spectroscopic velocity dispersions from DESI observations \citep{2025MNRAS.539.3627S}. The FP zero-point calibration was established using a collection of galaxies in the Coma cluster, with the distance of the latter constrained by 13 Type Ia supernovae within Coma analyzed by \citet{2025ApJ...979L...9S}, as well as an additional (fixed) constraint from SBF; see Appendix~\ref{sec:eq_special}.
%\riacomm{What does `additional calibration' mean in this case? App.\ref{sec:eq_special:FP} suggests that SNe\,Ia alone set the absolute scale?}

\paragraph{Tully-Fisher (TF)\label{sec:data:TF}}
The Tully-Fisher relation relates the rotation velocity of spiral galaxies (determined with HI line widths) to the total intrinsic luminosity. TF data were obtained from the Cosmicflows-4 catalog \citep{2020ApJ...902..145K}, which compiled HI line widths, redshifts, and photometry for $\sim 10,000$ spiral galaxies across the full sky, out of which we use 3400 galaxies with complete infrared photometry (as recommended by P.~Boubel 2025, priv.\ comm.). The TF relation zero-point has been re-calibrated within the {\DN} using calibrator galaxies identified in \citet{2024MNRAS.533.1550B}, with systematic corrections applied by \citet{2024arXiv241208449S}.
